\chapter{Conclusiones y líneas futuras}

\section{Conclusiones}
\par
Este trabajo ha permitido desarrollar e integrar un sistema completo de navegaci\'on para un UAV en entornos interiores, 
combinando percepci\'on, estimaci\'on, planificaci\'on y control dentro de una misma arquitectura funcional. El resultado 
final no debe entenderse como una suma de bloques aislados, sino como una soluci\'on conjunta en la que cada m\'odulo aporta 
una parte necesaria para conseguir un comportamiento global coherente.

Desde el punto de vista de la arquitectura del sistema, se ha definido una separaci\'on adecuada entre la plataforma a\'erea 
y la estaci\'on de tierra, lo que ha permitido aprovechar mejor los recursos de c\'omputo disponibles y simplificar el 
desarrollo experimental. De forma complementaria, el sistema de comunicaciones implementado ha permitido transportar de 
manera estable la informaci\'on visual, inercial y de control, resolviendo de forma pr\'actica las necesidades del proyecto.

En la parte de estimaci\'on, el sistema de odometr\'ia visual-inercial desarrollado ha demostrado ser capaz de proporcionar 
una estimaci\'on de pose utilizable para navegaci\'on. Los ensayos realizados han mostrado tanto las ventajas de la fusi\'on 
visual-inercial en el sistema completo como las limitaciones de una soluci\'on puramente visual, especialmente en escala y 
en maniobras con giro. Esto confirma que la integraci\'on de la IMU aporta una mejora real en robustez y consistencia.

En paralelo, el estimador de profundidad y el sistema de evitaci\'on de obst\'aculos han permitido dotar al dron de una 
capacidad reactiva frente a obst\'aculos cercanos. Los resultados obtenidos muestran que la direcci\'on de evitaci\'on y los 
comandos generados son coherentes con la escena observada, lo que valida el enfoque adoptado para esta parte del problema.

Por otro lado, la capa de planificaci\'on y control ha permitido transformar una trayectoria global en referencias locales y, 
a partir de ellas, en comandos de velocidad ejecutables por el veh\'iculo. La validaci\'on experimental ha mostrado que el 
planificador local selecciona puntos de seguimiento coherentes con la trayectoria y que esta decisi\'on se refleja 
correctamente en los comandos de gui\~nada y movimiento.

Finalmente, el conjunto de ensayos realizados ha permitido validar de forma progresiva los distintos bloques del sistema, 
desde pruebas sint\'eticas de preintegraci\'on IMU hasta evaluaciones del sistema completo. En conjunto, puede concluirse 
que se han cumplido los objetivos principales del proyecto: disponer de una base funcional de navegaci\'on aut\'onoma para un 
UAV, integrar los distintos subsistemas en una \'unica cadena de funcionamiento y demostrar experimentalmente que el sistema 
responde de forma coherente en los escenarios evaluados.
No obstante, los ensayos temporales tambi\'en muestran que la frecuencia global de funcionamiento todav\'ia no es tan alta 
como ser\'ia deseable, especialmente por la carga computacional asociada al bloque VIO, por lo que el sistema a\'un presenta 
un margen claro de mejora en rendimiento.


\section{Líneas futuras}

A partir de los resultados obtenidos en este trabajo, se identifican varias líneas de mejora que podrían abordarse en desarrollos futuros. 
Aunque la solución propuesta permite avanzar hacia una navegación segura de un UAV en entornos interiores, todavía presenta ciertas 
limitaciones. Por ello, a continuación se recogen algunos puntos de mejora que, por cuestiones de tiempo o por quedar fuera del alcance 
definido para este proyecto, no se han desarrollado en profundidad, pero que podrían aportar un valor significativo a la solución propuesta:

\begin{itemize}
    \item Optimizar computacionalmente el sistema, en especial el bloque de odometr\'ia visual-inercial y el c\'alculo de referencia, para aumentar la frecuencia global de funcionamiento y acercar la soluci\'on a un comportamiento m\'as pr\'oximo a tiempo real.
    \item Desarrollar un sistema VI-SLAM completo, incorporando la generación de un mapa del entorno y técnicas de \textit{loop closure}.
    \item Implementar un planificador global que tenga en cuenta el mapa del entorno y genere una trayectoria en función de la construcción del mapa en línea.
    \item Aumentar el tiempo de calibración mediante varianza de Allan hasta, al menos, 24 a 48 horas, con el objetivo de obtener una caracterización más precisa de la IMU.
    \item Implementar una interfaz gráfica de control que permita utilizar el algoritmo de forma más visual e intuitiva.
    \item Explorar e investigar algoritmos de planificación más avanzados para UAV en entornos interiores.
    \item Validación experimental en escenarios interiores más complejos.
\end{itemize}
